\section{Missing transverse momentum}
\label{sec:ch6_met}

\subsection{Reconstruction}

The incoming partons carry little transverse momentum compared with the hard-scattering scale, so an imbalance in the measured transverse momentum can indicate undetected particles.
Both the neutrino from the hard process and the tau-decay neutrino contribute to the signal imbalance.
The reconstructed vector is obtained from the negative sum of calibrated visible contributions~\cite{ATLAS2025METPerformance}:
\begin{align}
    \vec p_{\mathrm{T}}^{\,\mathrm{miss}}
    &=-\left(
    \sum_e\vec p_{\mathrm{T}}^{\,e}
    +\sum_\mu\vec p_{\mathrm{T}}^{\,\mu}
    +\sum_{\tau_{\mathrm{had}}}\vec p_{\mathrm{T}}^{\,\tau}
    +\sum_j\vec p_{\mathrm{T}}^{\,j}
    +\vec p_{\mathrm{T}}^{\,\mathrm{soft}}
    \right),
    \label{eq:ch6_met_vector}\\
    \met&=\left|\vec p_{\mathrm{T}}^{\,\mathrm{miss}}\right|.
    \label{eq:ch6_met_magnitude}
\end{align}
Two types of contribution are included:
\begin{itemize}
    \item \textbf{Hard term}: the electron, muon, hadronic-tau and jet contributions accepted by the reconstruction. No separate photon collection is supplied in this analysis.
    \item \textbf{Track soft term}: the vector sum of selected tracks associated with the primary vertex and not assigned to retained hard-object contributions.
\end{itemize}
The primary-vertex requirement suppresses charged-particle contributions from pile-up.
Neutral particles outside the hard term are not measured by the track soft term.

The baseline electron, muon and tau identification and isolation working points are used to select the corresponding inputs.
Jets are selected separately by the $\met$ reconstruction with the \texttt{Tight} operating point.
This requires jet $\pT>20~\GeV$ for $|\eta|<2.5$ and $\pT>30~\GeV$ for $2.5\leq|\eta|<4.5$.
The higher forward threshold suppresses pile-up jets beyond the tracking acceptance~\cite{ATLAS2025METPerformance}.
Their acceptance is distinct from that of the central jets counted in the analysis regions.
An association map is used to record shared detector signals and resolve double counting when the object contributions are combined~\cite{ATLAS2025METPerformance}.
Consequently, $\met$ is not recomputed as a simple sum of the objects passing the final region selections or the geometrical overlap-removal procedure.

\subsection{Calibration and uncertainties}

The energy and momentum calibrations of the hard objects are propagated to the reconstructed vector.
The response and resolution uncertainties of the track soft term are evaluated separately using events in which little genuine $\met$ is expected~\cite{ATLAS2025METPerformance}.
The corresponding variations are propagated together with the hard-object uncertainties to retain their correlations with the visible kinematics.

Apparent $\met$ can also be produced by mismeasured visible momentum, calorimeter noise or energy lost in uninstrumented regions.
For example, an underestimated jet momentum can produce an imbalance along the jet direction.
Jet cleaning and event-level angular selections are therefore used to suppress these contributions, as described in Chapter~\ref{chapter:event_selection}.
